\chapter{Results and Discussion}


This chapter follows the four-phase workflow established in Chapter 3. Phase I reports the available concrete design-\allowbreak{}check evidence for the representative Draft Final NBC 205:2012 regular building. Phase II establishes the NBC 205:2024 Regular model as the performance-based reference. Phase III evaluates the NBC 205:2024 plan-irregular buildings and distinguishes them from the separate four-storey height-extension model. Phase IV synthesizes the response parameters to determine practical RUD applicability limits.

The results presented in this chapter include NBC 105:2025 storey shear, storey drift, CM--CR separation, pushover response, performance points, hinge observations, concrete design-check observations for the representative NBC 205:2012 building, and the separate four-storey extension. All seven models in the main regular/plan-irregular comparison have three full storeys plus a stair cover. Only the separate height-extension model has four full storeys plus a stair cover, and it is not included in the seven-model performance-point tables.

\section{Phase I: Review of the Representative NBC 205:2012 Building}
\subsection{Evaluation of the 2012 Reference Building}
The representative building originally designed using Draft Final NBC 205:2012 was evaluated under five NBC 105:2025 seismic demand cases: Soil Type C with $Z=0.25$, $Z=0.30$, $Z=0.35$ and $Z=0.40$, and Soil Type D with $Z=0.35$. The concrete frame design-check results are compared below.

\subsection{Concrete Design-Check Observations}
\begin{table}[!htbp]
\caption[Concrete design-check failures for the 2012 reference]{Concrete frame design-check failures observed for the representative NBC 205:2012 building.}
\label{tab:ch4_2012_failures}
\centering
\begin{tabular}{@{}ccc@{}}
\toprule
\textbf{Soil} & \textbf{Z} & \textbf{Frames reported as failing the concrete design check}\\
\midrule
C & 0.25 & 1\\
C & 0.30 & 15\\
C & 0.35 & 21\\
C & 0.40 & 21\\
D & 0.35 & 18\\
\bottomrule
\end{tabular}

\vspace{0.5ex}\par\noindent\footnotesize\textit{Note: The counts identify deficient frame members in the investigated model and are not a general failure rate for all buildings designed under NBC 205:2012.}
\end{table}

The tabulated concrete design-\allowbreak{}check summary shows that the number of frame members failing the concrete design check increased substantially for the higher-demand cases. The largest reported count was 21 frames for Soil Type C at both $Z=0.35$ and $Z=0.40$, followed by 18 frames for Soil Type D at $Z=0.35$, 15 frames for Soil Type C at $Z=0.30$, and 1 frame for Soil Type C at $Z=0.25$. These results describe the adequacy of the investigated NBC 205:2012-designed building under the NBC 105:2025 demand cases considered in this study.

\section{Phase II: NBC 205:2024 Regular Benchmark}
The NBC 205:2024 Regular building is the reference model for the performance-based comparisons in the remainder of this chapter. It has three full storeys plus a stair cover and is evaluated under the same five NBC 105:2025 soil/seismic demand cases used for the seven-model comparison.

The Regular model therefore provides the comparison basis for the L, L1, L2, T, U, and Non-Parallel configurations. Performance-point base shear is interpreted as a demand--capacity intersection and is kept distinct from peak pushover capacity.

\section{Phase III: Performance Evaluation of Irregular NBC 205:2024 Buildings}

\subsection{Modal Characteristics}
Table~\ref{tab:modal_period_mass_z040} compares the first three natural periods and the dominant translational modal mass participation of the seven building configurations under Soil Type C at $Z=0.40$. This common demand case permits comparison of their dynamic characteristics.

\begin{table}[!htbp]
\centering
\caption{First three modal periods and dominant translational mass participation for Soil Type C, $Z=0.40$.}
\label{tab:modal_period_mass_z040}
\small
\begin{tabular}{@{}lccccc@{}}
\toprule
Model & $T_1$ (s) & $T_2$ (s) & $T_3$ (s) & Mode 1 UX (\%) & Mode 2 UY (\%) \\
\midrule
Regular & 0.699 & 0.661 & 0.569 & 80.32 & 80.12 \\
L & 0.698 & 0.660 & 0.564 & 79.95 & 80.42 \\
L1 & 0.685 & 0.649 & 0.561 & 79.14 & 80.16 \\
L2 & 0.649 & 0.593 & 0.476 & 79.83 & 79.07 \\
T & 0.698 & 0.674 & 0.569 & 79.09 & 74.82 \\
U & 0.707 & 0.649 & 0.569 & 79.30 & 78.77 \\
Non-Parallel & 0.650 & 0.599 & 0.489 & 79.14 & 78.75 \\
\bottomrule
\end{tabular}
\end{table}


The first mode is predominantly X-translational for every configuration, contributing 79.09--80.32\% of the translational mass in UX. The second mode is predominantly Y-translational, contributing 74.82--80.42\% in UY. The fundamental period $T_1$ ranges from 0.649~s for L2 to 0.707~s for U; $T_2$ ranges from 0.593 to 0.674~s, while $T_3$ ranges from 0.476 to 0.569~s. These differences indicate modest but measurable changes in the dynamic characteristics among the investigated plan configurations and do not by themselves establish a performance ranking.

\begin{table}[!htbp]
\centering
\caption{Cumulative translational modal mass participation for Soil Type C, $Z=0.40$.}
\label{tab:modal_cumulative_mass_z040}
\small
\resizebox{\textwidth}{!}{%
\begin{tabular}{@{}lcccccc@{}}
\toprule
Model & SumUX at Mode 3 (\%) & SumUY at Mode 3 (\%) & X $\geq$ 90\% mode & SumUX (\%) & Y $\geq$ 90\% mode & SumUY (\%) \\
\midrule
Regular & 80.67 & 81.86 & 4 & 92.39 & 5 & 91.89 \\
L & 80.18 & 81.35 & 4 & 92.25 & 5 & 91.96 \\
L1 & 79.38 & 80.56 & 4 & 91.94 & 5 & 91.71 \\
L2 & 79.84 & 81.49 & 4 & 91.96 & 5 & 92.03 \\
T & 79.51 & 80.75 & 4 & 92.04 & 5 & 90.94 \\
U & 79.94 & 81.60 & 4 & 92.01 & 5 & 91.38 \\
Non-Parallel & 80.44 & 81.90 & 4 & 92.31 & 5 & 92.39 \\
\bottomrule
\end{tabular}
}
\end{table}


For every configuration, cumulative translational participation exceeds 90\% by Mode~4 in the X direction and by Mode~5 in the Y direction. The corresponding cumulative values at the threshold modes are approximately 91.94--92.39\% in X and 90.94--92.39\% in Y. By Mode~12, cumulative translational mass participation reaches 100\% in both directions for all seven configurations. The modal results therefore provide an explicit direction-wise participation check in addition to the period comparison used elsewhere in the response discussion.
\FloatBarrier


\subsection{Building Configurations Considered in the Study}
The building configurations used in the analysis were introduced in Chapter 3 and are illustrated in Figures~\ref{fig:method_regular_model} to~\ref{fig:method_np_model}. The Regular building was used as the reference model, while the L-shaped, L1-shaped, L2-shaped, T-shaped, and U-shaped buildings represent common re-entrant plan irregularities. The Non-Parallel building represents a non-orthogonal framing arrangement with skewed load paths.

The comparison among these seven models is important because NBC 205 ready-to-use provisions are primarily intended for ordinary low-rise buildings, whereas practical residential layouts often depart from regular rectangular plans. Plan irregularity can alter stiffness distribution, mass distribution, centre of rigidity location, and lateral force transfer. These effects may not be fully understood from strength checks alone; therefore, the following sections compare force demand, drift, descriptive CM--CR separation, capacity, performance-point characteristics, and observed hinge/failure evidence together.

The seven building configurations have three full storeys plus a stair cover, allowing their responses to be compared at the same building height. The separate four-storey extension is discussed through its own design-check results.

\subsection{Storey Shear Behaviour}
Storey shear is a primary indicator of lateral seismic force demand. Figures~\ref{fig:story_shear_group1} and~\ref{fig:story_shear_group2} show the storey-shear response of the seven building configurations over the investigated demand cases.

\begin{figure}[!htbp]
\centering
\begin{subfigure}[t]{0.48\textwidth}\centering
\includegraphics[width=\linewidth]{figures/ch4/cp114_panels/storey_shear_regular.pdf}\subcaption{Regular}
\end{subfigure}\hfill
\begin{subfigure}[t]{0.48\textwidth}\centering
\includegraphics[width=\linewidth]{figures/ch4/cp114_panels/storey_shear_l.pdf}\subcaption{L-shaped}
\end{subfigure}
\begin{subfigure}[t]{0.48\textwidth}\centering
\includegraphics[width=\linewidth]{figures/ch4/cp114_panels/storey_shear_l1.pdf}\subcaption{L1-shaped}
\end{subfigure}\hfill
\begin{subfigure}[t]{0.48\textwidth}\centering
\includegraphics[width=\linewidth]{figures/ch4/cp114_panels/storey_shear_l2.pdf}\subcaption{L2-shaped}
\end{subfigure}
\caption{Storey shear: Regular, L, L1 and L2.}
\label{fig:story_shear_group1}
\end{figure}

\begin{figure}[!htbp]
\centering
\begin{subfigure}[t]{0.48\textwidth}\centering
\includegraphics[width=\linewidth]{figures/ch4/cp114_panels/storey_shear_t.pdf}\subcaption{T-shaped}
\end{subfigure}\hfill
\begin{subfigure}[t]{0.48\textwidth}\centering
\includegraphics[width=\linewidth]{figures/ch4/cp114_panels/storey_shear_u.pdf}\subcaption{U-shaped}
\end{subfigure}
\begin{subfigure}[t]{0.48\textwidth}\centering
\includegraphics[width=\linewidth]{figures/ch4/cp114_panels/storey_shear_np.pdf}\subcaption{Non-Parallel}
\end{subfigure}
\caption{Storey shear: T, U and Non-Parallel.}
\label{fig:story_shear_group2}
\end{figure}
\FloatBarrier

The storey-shear results show that the largest value occurred in the Regular building under Soil Type C and seismic zone factor $Z=0.40$. The Story 1 shear was approximately 532.489~kN. Under the same condition and storey level, the U-shaped, Non-Parallel, and L-shaped buildings reached approximately 504.048~kN, 486.757~kN, and 476.287~kN, respectively.

The comparison does not show a universal increase in storey shear simply because a plan is irregular. The Regular model governs the maximum value, while several irregular models remain comparatively close to it. Plan geometry therefore changes the distribution of force demand, but its effect must be interpreted model by model rather than assumed to be an automatic amplification.

The maximum storey shear values for most models occurred at Story 1, which is consistent with the accumulation of lateral inertia forces toward the base. Across the storey shear results, the controlling values generally occurred under Soil Type C with $Z=0.40$, indicating that the highest seismic zone factor governed the largest force demand. Soil Type D at $Z=0.35$ did not govern the maximum storey shear among the analysed cases.

From the viewpoint of ready-to-use seismic adequacy, the storey-shear comparison still shows the need to examine the analysed force distribution rather than infer demand from plan shape alone. The Regular model governed the lower-storey shear, while the irregular models produced different but substantial force profiles under the same seismic design framework.


\begin{table}[!htbp]
\centering
\caption{Maximum storey shear for Soil Type C, $Z=0.40$.}
\label{tab:storey_shear_summary_cp112}
\small
\resizebox{\textwidth}{!}{%
\begin{tabular}{@{}lcc@{}}
\toprule
Model & Max. EQx ULS storey shear (kN) & Governing available case \\
\midrule
Regular & 532.489 & EQx, Story 1 \\
L & 476.287 & EQx, Story 1 \\
L1 & 403.859 & EQx, Story 1 \\
L2 & 456.660 & EQx, Story 1 \\
T & 407.396 & EQx, Story 1 \\
U & 504.048 & EQx, Story 1 \\
Non-Parallel & 486.757 & EQx, Story 1 \\
\bottomrule
\end{tabular}
}
\par\smallskip
\begin{minipage}{0.97\textwidth}\footnotesize
The available storey-shear source supports the EQx comparison shown here. The Regular Soil C, $Z=0.30$ force-scale anomaly identified during data checking is not used in this table.
\end{minipage}
\end{table}


\subsection{Centre of Mass and Centre of Rigidity Separation}
CM--CR separation is reported descriptively and is considered together with the other response measures. Figure~\ref{fig:nbc105_2025_cmcr_main_story} and Table~\ref{tab:ch4_cmcr_summary} summarize the maximum resultant separation across main Stories 1--3. Figures~\ref{fig:cmcr_soil_c} and~\ref{fig:cmcr_soil_d} retain the detailed comparison for Soil Types C and D.

\begin{figure}[!htbp]
\centering
\includegraphics[width=0.92\textwidth]{figures/ch4/authority_exact/fig_4_8_CP110_exact_clip.pdf}
\caption{CM--CR separation for Soil Type C.}
\label{fig:cmcr_soil_c}
\end{figure}

\begin{figure}[!htbp]
\centering
\includegraphics[width=0.92\textwidth]{figures/ch4/authority_exact/fig_4_9_CP110_exact_clip.pdf}
\caption{CM--CR separation for Soil Type D.}
\label{fig:cmcr_soil_d}
\end{figure}
\FloatBarrier

The largest main-storey resultant CM--CR separation in the analysed dataset occurs in the Non-Parallel building under Soil Type C and $Z=0.25$, reaching approximately 1.237~m at Story 2. The L2-shaped building follows at approximately 1.186~m under Soil Type C and $Z=0.25$ at Story 1. The corresponding model maxima for Regular, L, L1, T, and U are approximately 0.398, 0.581, 0.640, 0.475, and 0.424~m, respectively.

The magnitude and governing demand case vary by plan configuration. CM--CR separation describes mass--stiffness eccentricity; it is not the diaphragm-edge displacement ratio used for code-based torsional-irregularity classification. It is therefore interpreted alongside storey shear, drift, pushover capacity, performance-point response and hinge-formation patterns.

\subsection{Torsional Behaviour}
Torsional response is important in plan-irregular buildings because eccentricity between mass and stiffness and non-orthogonal framing can induce floor rotation. Figures~\ref{fig:torsion_c025} to~\ref{fig:torsion_d035} present the torsional rotation comparisons retained from the approved analysis.

\begin{figure}[!htbp]
\centering
\includegraphics[width=0.90\textwidth]{figures/ch4/authority_exact/fig_4_10_CP110_exact_clip.pdf}
\caption{Torsional rotation: Soil C, $Z=0.25$.}
\label{fig:torsion_c025}
\end{figure}
\begin{figure}[!htbp]
\centering
\includegraphics[width=0.90\textwidth]{figures/ch4/authority_exact/fig_4_11_CP110_exact_clip.pdf}
\caption{Torsional rotation: Soil C, $Z=0.30$.}
\label{fig:torsion_c030}
\end{figure}
\begin{figure}[!htbp]
\centering
\includegraphics[width=0.90\textwidth]{figures/ch4/authority_exact/fig_4_12_CP110_exact_clip.pdf}
\caption{Torsional rotation: Soil C, $Z=0.35$.}
\label{fig:torsion_c035}
\end{figure}
\begin{figure}[!htbp]
\centering
\includegraphics[width=0.90\textwidth]{figures/ch4/authority_exact/fig_4_13_CP110_exact_clip.pdf}
\caption{Torsional rotation: Soil C, $Z=0.40$.}
\label{fig:torsion_c040}
\end{figure}
\begin{figure}[!htbp]
\centering
\includegraphics[width=0.90\textwidth]{figures/ch4/authority_exact/fig_4_14_CP110_exact_clip.pdf}
\caption{Torsional rotation: Soil D, $Z=0.35$.}
\label{fig:torsion_d035}
\end{figure}
\FloatBarrier

\subsection{Nonlinear Static Pushover Capacity Behaviour}
Nonlinear static pushover analysis was used to evaluate the lateral load capacity and deformation characteristics of the seven building configurations. Figures~\ref{fig:cap_regular_x} to~\ref{fig:cap_np_y} present the capacity curves of the seven building configurations in the X and Y directions. The curves are presented through the peak base-shear point for comparative presentation, while performance-point demand is discussed separately.

For consistent comparison among the building configurations, the capacity curves were presented up to the peak base-shear point. Subsequent descending and displacement-\allowbreak{}reversing points associated with post-peak degradation or terminal numerical behaviour were excluded from the comparative figures. The original analysis outputs were retained for verification.

The capacity-curve figures are interpreted in terms of the pre-peak curve shape, relative lateral stiffness, peak base shear, and the displacement corresponding to the first global peak. Because the figures stop at that first peak, they are not used here to infer post-peak ductility, residual strength, collapse displacement, or post-peak degradation rate.

The highest first-peak pushover base shear in the capacity results was obtained for the U-shaped building under Soil Type C and $Z=0.40$ in the Y direction. The peak base shear was approximately 1882.070~kN at a monitored displacement of approximately 174.697~mm. For the Regular building under the same soil condition, seismic zone factor, and Y-direction pushover, the corresponding peak base shear was approximately 1807.405~kN at 147.246~mm.

The U-shaped building therefore showed slightly higher ultimate pushover strength than the Regular model in the critical capacity case. However, this should not be interpreted as meaning that the U-shaped building has the best overall seismic performance. Ultimate capacity represents only one part of seismic behaviour. The same U-shaped model also governed the performance-point roof displacement, while the Non-Parallel model governed the main-storey CM--CR separation comparison.

The capacity results also show that several irregular configurations developed substantial lateral strength before reaching their respective peak base-shear points. For example, the Non-Parallel building reached a high ultimate pushover base shear capacity in its governing pushover direction, while the L2-shaped building also showed a comparatively high capacity. These results indicate that NBC 205:2024 irregular buildings may possess significant strength reserve, but strength reserve alone is not sufficient for judging seismic adequacy.

The comparison among the seven models and five soil/seismic cases shows that plan geometry affects both the slope of the pre-peak response and the peak base-shear level. The Regular model provides the reference response, while the L-, L1-, L2-, T-, U-, and Non-Parallel models show how re-entrant corners and non-orthogonal framing can change stiffness distribution and the displacement required to reach peak resistance. For performance-based evaluation, these capacity-curve characteristics are interpreted together with separate performance-\allowbreak{}point demand, drift, CM--CR separation, and the available hinge/failure evidence.


\begin{figure}[!htbp]
\centering
\includegraphics[width=0.90\textwidth,trim=0 0 0 24bp,clip]{figures/ch4/authority_exact/fig_4_15_CP97_CP66_exact_clip.pdf}
\caption{Regular building capacity curve: X direction.}
\label{fig:cap_regular_x}
\end{figure}

\begin{figure}[!htbp]
\centering
\includegraphics[width=0.90\textwidth,trim=0 0 0 24bp,clip]{figures/ch4/authority_exact/fig_4_16_CP97_CP66_exact_clip.pdf}
\caption{Regular building capacity curve: Y direction.}
\label{fig:cap_regular_y}
\end{figure}

\begin{figure}[!htbp]
\centering
\includegraphics[width=0.90\textwidth,trim=0 0 0 24bp,clip]{figures/ch4/authority_exact/fig_4_17_CP97_CP66_exact_clip.pdf}
\caption{L-shaped building capacity curve: X direction.}
\label{fig:cap_l_x}
\end{figure}

\begin{figure}[!htbp]
\centering
\includegraphics[width=0.90\textwidth,trim=0 0 0 24bp,clip]{figures/ch4/authority_exact/fig_4_18_CP97_CP66_exact_clip.pdf}
\caption{L-shaped building capacity curve: Y direction.}
\label{fig:cap_l_y}
\end{figure}

\begin{figure}[!htbp]
\centering
\includegraphics[width=0.90\textwidth,trim=0 0 0 24bp,clip]{figures/ch4/authority_exact/fig_4_19_CP97_CP66_exact_clip.pdf}
\caption{L1-shaped building capacity curve: X direction.}
\label{fig:cap_l1_x}
\end{figure}

\begin{figure}[!htbp]
\centering
\includegraphics[width=0.90\textwidth,trim=0 0 0 24bp,clip]{figures/ch4/authority_exact/fig_4_20_CP97_CP66_exact_clip.pdf}
\caption{L1-shaped building capacity curve: Y direction.}
\label{fig:cap_l1_y}
\end{figure}

\begin{figure}[!htbp]
\centering
\includegraphics[width=0.90\textwidth,trim=0 0 0 24bp,clip]{figures/ch4/authority_exact/fig_4_21_CP97_CP66_exact_clip.pdf}
\caption{L2-shaped building capacity curve: X direction.}
\label{fig:cap_l2_x}
\end{figure}

\begin{figure}[!htbp]
\centering
\includegraphics[width=0.90\textwidth,trim=0 0 0 24bp,clip]{figures/ch4/authority_exact/fig_4_22_CP97_CP66_exact_clip.pdf}
\caption{L2-shaped building capacity curve: Y direction.}
\label{fig:cap_l2_y}
\end{figure}

\begin{figure}[!htbp]
\centering
\includegraphics[width=0.90\textwidth,trim=0 0 0 24bp,clip]{figures/ch4/authority_exact/fig_4_23_CP97_CP66_exact_clip.pdf}
\caption{T-shaped building capacity curve: X direction.}
\label{fig:cap_t_x}
\end{figure}

\begin{figure}[!htbp]
\centering
\includegraphics[width=0.90\textwidth,trim=0 0 0 24bp,clip]{figures/ch4/authority_exact/fig_4_24_CP97_CP66_exact_clip.pdf}
\caption{T-shaped building capacity curve: Y direction.}
\label{fig:cap_t_y}
\end{figure}

\begin{figure}[!htbp]
\centering
\includegraphics[width=0.90\textwidth,trim=0 0 0 24bp,clip]{figures/ch4/authority_exact/fig_4_25_CP97_CP66_exact_clip.pdf}
\caption{U-shaped building capacity curve: X direction.}
\label{fig:cap_u_x}
\end{figure}

\begin{figure}[!htbp]
\centering
\includegraphics[width=0.90\textwidth,trim=0 0 0 24bp,clip]{figures/ch4/authority_exact/fig_4_26_CP97_CP66_exact_clip.pdf}
\caption{U-shaped building capacity curve: Y direction.}
\label{fig:cap_u_y}
\end{figure}

\begin{figure}[!htbp]
\centering
\includegraphics[width=0.90\textwidth,trim=0 0 0 24bp,clip]{figures/ch4/authority_exact/fig_4_27_CP97_CP66_exact_clip.pdf}
\caption{Non-Parallel building capacity curve: X direction.}
\label{fig:cap_np_x}
\end{figure}

\begin{figure}[!htbp]
\centering
\includegraphics[width=0.90\textwidth,trim=0 0 0 24bp,clip]{figures/ch4/authority_exact/fig_4_28_CP97_CP66_exact_clip.pdf}
\caption{Non-Parallel building capacity curve: Y direction.}
\label{fig:cap_np_y}
\end{figure}
\FloatBarrier


\begin{table}[!htbp]
\centering
\caption{Seven-model pushover and performance-point summary.}
\label{tab:pushover_pp_summary_cp112}
\scriptsize
\resizebox{\textwidth}{!}{%
\begin{tabular}{@{}lcrrcrr@{}}
\toprule
Model & Peak dir. & Peak shear & Peak disp. & PP dir. & Max PP shear & Max PP disp. \\
 & & (kN) & (mm) & & (kN) & (mm) \\
\midrule
Regular & X & 1836.338 & 198.040 & X & 456.253 & 26.053 \\
L & Y & 1800.120 & 187.539 & X & 407.295 & 25.900 \\
L1 & X & 1606.250 & 212.287 & X & 340.753 & 24.592 \\
L2 & Y & 1804.152 & 172.789 & X & 370.935 & 21.171 \\
T & Y & 1640.715 & 223.937 & X & 343.488 & 24.802 \\
U & Y & 1882.070 & 174.697 & X & 435.084 & 27.072 \\
Non-Parallel & Y & 1827.694 & 154.530 & X & 399.636 & 21.328 \\
\bottomrule
\end{tabular}
}
\par\smallskip
\begin{minipage}{0.97\textwidth}\footnotesize
Table~\ref{tab:pushover_pp_summary_cp112} shows the peak pushover base shear and displacement for each building configuration, alongside the maximum reported performance-point values. The peak values describe the first global capacity peak on the recorded curves, whereas the performance-point quantities describe the available demand-state intersection. For all seven models, the reported maximum performance-point base shear and roof displacement occur in the X direction; these values should not be read as the corresponding peak-pushover ordinates.
\end{minipage}
\end{table}


\subsection{Performance Point Characteristics}
The performance point identifies the estimated seismic demand state from the pushover analysis. It is therefore different from the ultimate pushover base shear capacity reported from the peak of the capacity curve. In this section, performance-\allowbreak{}point base shear and performance-\allowbreak{}point roof displacement are presented in tables only, so that the demand quantities remain distinct from the peak capacity quantities shown in the capacity-curve figures.

\begin{table}[!htbp]
\caption{Maximum performance-\allowbreak{}point base shear and roof displacement for the seven three-storey building configurations.}
\label{tab:ch4_pp_max}
\centering
\footnotesize
\setlength{\tabcolsep}{3pt}
\begin{tabular}{@{}l r c c c r c c c@{}}
\toprule
\textbf{Model} & \shortstack{\textbf{Max. PP base}\\\textbf{shear (kN)}} & \textbf{Soil} & \textbf{Z} & \textbf{Dir.} & \shortstack{\textbf{Max. PP roof}\\\textbf{disp. (mm)}} & \textbf{Soil} & \textbf{Z} & \textbf{Dir.}\\
\midrule
Regular & 456.253 & C & 0.40 & X & 26.053 & C & 0.40 & X\\
L-shaped & 407.295 & C & 0.40 & X & 25.900 & C & 0.40 & X\\
L1-shaped & 340.753 & C & 0.40 & X & 24.592 & C & 0.40 & X\\
L2-shaped & 370.935 & C & 0.40 & X & 21.171 & C & 0.40 & X\\
T-shaped & 343.488 & C & 0.40 & X & 24.802 & C & 0.40 & X\\
U-shaped & 435.084 & C & 0.40 & X & 27.072 & C & 0.40 & X\\
Non-Parallel & 399.636 & C & 0.40 & X & 21.328 & C & 0.40 & X\\
\bottomrule
\end{tabular}
\end{table}

The performance-point results show that the highest performance-\allowbreak{}point base shear was obtained for the Regular building under Soil Type C and $Z=0.40$ in the X direction, with a value of approximately 456.253~kN and a corresponding roof displacement of approximately 26.053~mm.

In contrast, the largest performance-\allowbreak{}point roof displacement was obtained for the U-shaped building under Soil Type C and $Z=0.40$ in the X direction. The displacement was approximately 27.072~mm, with a corresponding performance-\allowbreak{}point base shear of approximately 435.084~kN. Under the same condition and direction, the Regular building had a performance-\allowbreak{}point roof displacement of approximately 26.053~mm and a performance-\allowbreak{}point base shear of approximately 456.253~kN.

These findings distinguish strength capacity from seismic demand at the performance point. The U-shaped model had the highest ultimate pushover base shear capacity, but the Regular model had the highest performance-\allowbreak{}point base shear. Similarly, the U-shaped model governed the performance-\allowbreak{}point roof displacement. Therefore, no single response parameter is sufficient to identify the best-performing model.

For applicability assessment, the performance-point comparison shows that irregular configurations should not be compared only by ultimate strength; the displacement and base-shear demand at the embedded performance point provide separate response measures. Formal adequacy still requires the project-specific checks applicable to the actual irregular building.

\newcommand{\pptable}[3]{%
\begin{table}[!htbp]
\caption{Performance points for the #1 building.}
\label{#2}
\centering
\footnotesize
\setlength{\tabcolsep}{4pt}
\begin{tabular}{@{}ccrrrr@{}}
\toprule
 & & \multicolumn{2}{c}{\textbf{X direction}} & \multicolumn{2}{c}{\textbf{Y direction}}\\
\textbf{Soil} & \textbf{Z} & \shortstack{\textbf{Base shear}\\\textbf{(kN)}} & \shortstack{\textbf{Roof displacement}\\\textbf{(mm)}} & \shortstack{\textbf{Base shear}\\\textbf{(kN)}} & \shortstack{\textbf{Roof displacement}\\\textbf{(mm)}}\\
\midrule
#3
\bottomrule
\end{tabular}
\end{table}}

\pptable{Regular}{tab:pp_regular}{C & 0.25 & 272.662 & 18.103 & 260.818 & 16.475\\ C & 0.30 & 326.721 & 20.736 & 312.562 & 18.884\\ C & 0.35 & 384.293 & 22.944 & 365.801 & 19.544\\ C & 0.40 & 456.253 & 26.053 & 425.945 & 21.748\\ D & 0.35 & 344.262 & 21.849 & 331.579 & 20.033\\}
\pptable{L-shaped}{tab:pp_l}{C & 0.25 & 227.469 & 16.766 & 221.517 & 15.837\\ C & 0.30 & 291.339 & 20.732 & 278.035 & 18.560\\ C & 0.35 & 343.154 & 22.823 & 324.218 & 19.079\\ C & 0.40 & 407.295 & 25.900 & 377.501 & 21.231\\ D & 0.35 & 307.087 & 21.852 & 295.011 & 19.694\\}
\pptable{L1-shaped}{tab:pp_l1}{C & 0.25 & 191.508 & 16.459 & 183.993 & 15.092\\ C & 0.30 & 236.086 & 19.386 & 227.704 & 17.854\\ C & 0.35 & 288.192 & 21.757 & 273.017 & 18.091\\ C & 0.40 & 340.753 & 24.592 & 320.987 & 20.320\\ D & 0.35 & 260.104 & 21.152 & 249.116 & 18.816\\}
\pptable{L2-shaped}{tab:pp_l2}{C & 0.25 & 216.371 & 14.337 & 206.176 & 12.121\\ C & 0.30 & 256.303 & 16.484 & 242.133 & 13.651\\ C & 0.35 & 321.068 & 19.170 & 297.992 & 14.987\\ C & 0.40 & 370.935 & 21.171 & 348.199 & 16.701\\ D & 0.35 & 285.319 & 18.090 & 271.252 & 15.257\\}
\pptable{T-shaped}{tab:pp_t}{C & 0.25 & 203.430 & 17.352 & 202.163 & 17.516\\ C & 0.30 & 255.839 & 20.791 & 256.484 & 21.313\\ C & 0.35 & 302.142 & 22.824 & 290.652 & 21.025\\ C & 0.40 & 343.488 & 24.802 & 330.261 & 22.826\\ D & 0.35 & 269.803 & 21.925 & 273.922 & 22.762\\}
\pptable{U-shaped}{tab:pp_u}{C & 0.25 & 250.799 & 18.081 & 241.465 & 15.610\\ C & 0.30 & 306.616 & 21.126 & 297.762 & 18.402\\ C & 0.35 & 366.562 & 23.849 & 346.440 & 22.540\\ C & 0.40 & 435.084 & 27.072 & 398.733 & 20.852\\ D & 0.35 & 325.614 & 21.430 & 314.949 & 18.601\\}
\pptable{Non-Parallel}{tab:pp_np}{C & 0.25 & 232.837 & 14.382 & 222.705 & 12.491\\ C & 0.30 & 290.176 & 17.157 & 274.571 & 14.743\\ C & 0.35 & 345.477 & 19.291 & 320.994 & 15.410\\ C & 0.40 & 399.636 & 21.328 & 375.298 & 17.168\\ D & 0.35 & 305.705 & 18.075 & 292.346 & 15.697\\}


\begin{figure}[!htbp]
\centering
\begin{subfigure}[t]{0.49\textwidth}\centering
\includegraphics[width=\linewidth]{figures/ch4/cp112_generated/pp_base_shear.png}\subcaption{Base shear}
\end{subfigure}\hfill
\begin{subfigure}[t]{0.49\textwidth}\centering
\includegraphics[width=\linewidth]{figures/ch4/cp112_generated/pp_displacement.png}\subcaption{Roof displacement}
\end{subfigure}
\caption{Performance-point response comparison.}
\label{fig:pp_model_comparison_cp112}
\end{figure}

Figure~\ref{fig:pp_model_comparison_cp112} shows that the model governing performance-point base shear is not the same as the model governing performance-point roof displacement. The Regular model has the largest performance-point base shear, whereas the U-shaped model has the largest roof displacement, reinforcing the need to interpret force and deformation demand separately.

\FloatBarrier

\subsection{Hinge-Evidence Summary}
The detailed final-step hinge screenshots remain in Appendix~\ref{appx:hinge_formation_results}. All seven models have X- and Y-direction final-step images for the five matched demand cases. However, the available analysis records do not provide a member-level table identifying a unique ``critical hinge location'', beam/column classification, and final hinge state for every model/direction. Table~\ref{tab:hinge_evidence_cp112} therefore reports the evidence status rather than inventing those fields from screenshots.
\begin{table}[!htbp]
\centering
\caption{Hinge evidence available for the seven-model comparison.}
\label{tab:hinge_evidence_cp112}
\small
\resizebox{\textwidth}{!}{%
\begin{tabular}{@{}lccp{6.3cm}@{}}
\toprule
Model & X final-step image & Y final-step image & Member-level result status \\
\midrule
Regular & Available & Available & Critical member/location/state not available in tabulated results \\
L & Available & Available & Same limitation \\
L1 & Available & Available & Same limitation \\
L2 & Available & Available & Same limitation \\
T & Available & Available & Same limitation \\
U & Available & Available & Same limitation \\
Non-Parallel & Available & Available & Same limitation \\
\bottomrule
\end{tabular}
}
\end{table}


\subsection{Observed Failures}
The observed-failure evidence is separated into two distinct comparisons so that the historical NBC 205:2012 building is not mixed with the NBC 205:2024 height-\allowbreak{}extension study.

\textbf{Existing NBC 205:2012-designed building under NBC 105:2025 demand}

The representative existing building designed using Draft Final NBC 205:2012 was checked under the five NBC 105:2025 demand cases used in Phase I. Table~\ref{tab:ch4_2012_failures} summarizes the identified failing frame counts, while Figure~\ref{fig:failure_2012_final_final} presents the corresponding design-check evidence. The observed counts were 1, 15, 21, 21, and 18 for Soil C/$Z=0.25$, C/$Z=0.30$, C/$Z=0.35$, C/$Z=0.40$, and Soil D/$Z=0.35$, respectively.

\begin{figure}[!htbp]
\centering
\begin{subfigure}[t]{0.32\textwidth}\centering
\includegraphics[width=\linewidth]{figures/ch4/final_final_original/failures_2012/SRC-FF-IMG-018_C025.png}\subcaption{Soil C, $Z=0.25$}
\end{subfigure}\hfill
\begin{subfigure}[t]{0.32\textwidth}\centering
\includegraphics[width=\linewidth]{figures/ch4/final_final_original/failures_2012/SRC-FF-IMG-014_C030.png}\subcaption{Soil C, $Z=0.30$}
\end{subfigure}\hfill
\begin{subfigure}[t]{0.32\textwidth}\centering
\includegraphics[width=\linewidth]{figures/ch4/final_final_original/failures_2012/SRC-FF-IMG-015_C035.png}\subcaption{Soil C, $Z=0.35$}
\end{subfigure}
\par\vspace{0.6em}
\begin{subfigure}[t]{0.32\textwidth}\centering
\includegraphics[width=\linewidth]{figures/ch4/final_final_original/failures_2012/SRC-FF-IMG-016_C040.png}\subcaption{Soil C, $Z=0.40$}
\end{subfigure}\hfill
\begin{subfigure}[t]{0.32\textwidth}\centering
\includegraphics[width=\linewidth]{figures/ch4/final_final_original/failures_2012/SRC-FF-IMG-017_D035.png}\subcaption{Soil D, $Z=0.35$}
\end{subfigure}
\caption{NBC 205:2012 design-check observations.}
\label{fig:failure_2012_final_final}
\end{figure}

This comparison represents an older NBC 205:2012 design practice being rechecked using the current NBC 105:2025 seismic-demand framework. It therefore provides a practical historical comparison with the current NBC 205:2024 Regular benchmark at the completed design-\allowbreak{}check level. Direct 2012-versus-2024 pushover comparison is outside the present completed scope.

\textbf{NBC 205:2024 four-storey Regular extension}

The separate four-storey Regular extension was examined under the same five soil and seismic zone-factor combinations used in the three-storey model comparisons. Its design-check observations are discussed independently of the seven-model response comparisons.

\begin{figure}[!htbp]
\centering
\begin{subfigure}[t]{0.32\textwidth}\centering
\includegraphics[width=\linewidth]{figures/ch4/final_final_original/failures_four_storey/SRC-FF-IMG-009_C025.png}\subcaption{Soil C, $Z=0.25$}
\end{subfigure}\hfill
\begin{subfigure}[t]{0.32\textwidth}\centering
\includegraphics[width=\linewidth]{figures/ch4/final_final_original/failures_four_storey/SRC-FF-IMG-010_C030.png}\subcaption{Soil C, $Z=0.30$}
\end{subfigure}\hfill
\begin{subfigure}[t]{0.32\textwidth}\centering
\includegraphics[width=\linewidth]{figures/ch4/final_final_original/failures_four_storey/SRC-FF-IMG-011_C035.png}\subcaption{Soil C, $Z=0.35$}
\end{subfigure}
\par\vspace{0.6em}
\begin{subfigure}[t]{0.32\textwidth}\centering
\includegraphics[width=\linewidth]{figures/ch4/final_final_original/failures_four_storey/SRC-FF-IMG-012_C040.png}\subcaption{Soil C, $Z=0.40$}
\end{subfigure}\hfill
\begin{subfigure}[t]{0.32\textwidth}\centering
\includegraphics[width=\linewidth]{figures/ch4/final_final_original/failures_four_storey/SRC-FF-IMG-013_D035.png}\subcaption{Soil D, $Z=0.35$}
\end{subfigure}
\caption{Four-storey extension design-check observations.}
\label{fig:failure_fourstorey_final_final}
\end{figure}

Figure~\ref{fig:failure_fourstorey_final_final} presents the four-storey extension design-check results under Soil Type C with $Z=0.25$, $0.30$, $0.35$ and $0.40$, and Soil Type D with $Z=0.35$. The observations describe the investigated extension and do not establish a general height-limit rule.

The combined failure evidence therefore supports two different conclusions. First, the existing NBC 205:2012-designed building exhibits deficient frame members when rechecked under the selected NBC 105:2025 demand cases. Second, within the current NBC 205:2024 the separate four-storey Regular extension shows design-\allowbreak{}check failures. The four-storey case highlights the importance of project-specific engineered checking beyond the scope of prescriptive RUD detailing.

\subsection{Seven-Model Response Under NBC 105:2025}
This section presents the NBC 105:2025 response results used in the main comparison. The seven-model comparison includes the NBC 205:2024 regular, irregular (L, L1, L2, T, U), and non-parallel configurations, all modelled as three full storeys plus a stair cover. The separate four-storey height-extension model is not merged into these seven-model comparison tables. Detailed numerical files are retained separately as supporting electronic records and are not reproduced as raw-data appendices in the printed thesis.

Table~\ref{tab:ch4_governing} summarizes the controlling response quantities used in the comparison. The largest storey shear is obtained by the Regular building at Story 1 under Soil Type C and $Z=0.40$. The largest first-peak pushover capacity is obtained by the U-shaped building in the Y direction under the same soil and zone condition, while the largest performance-\allowbreak{}point base shear is obtained by the Regular building in the X direction. The largest performance-\allowbreak{}point roof displacement is obtained by the U-shaped building in the X direction. These quantities are kept separate because storey shear, peak pushover resistance, and performance-\allowbreak{}point demand describe different aspects of structural response.

\begin{table}[!htbp]
\caption[Governing response quantities for seven models]{Selected governing response quantities for the seven-model comparison.}
\label{tab:ch4_governing}
\centering
\begin{tabularx}{\textwidth}{@{}p{0.25\textwidth} p{0.18\textwidth} X@{}}
\toprule
\textbf{Response quantity} & \textbf{Governing value} & \textbf{Governing model and condition}\\
\midrule
Storey shear & 532.489 kN & Regular, Story 1, Soil C, $Z=0.40$\\
Peak pushover base shear & 1882.070 kN & U-shaped, Y direction, Soil C, $Z=0.40$; displacement=174.697 mm\\
Performance-point base shear & 456.253 kN & Regular, X direction, Soil C, $Z=0.40$\\
Performance-point roof displacement & 27.072 mm & U-shaped, X direction, Soil C, $Z=0.40$\\
ULS drift ratio & 0.004954 & U-shaped, X direction, Soil C, $Z=0.40$, Story 2\\
SLS drift ratio & 0.004744 & U-shaped, X direction, Soil C, $Z=0.40$, Story 2; 79.067\% of the 0.006 SLS limit\\
Main-storey CM--CR separation & 1.237 m & Non-Parallel, Soil C, $Z=0.25$, Story 2\\
\bottomrule
\end{tabularx}
\end{table}

\begin{table}[!htbp]
\caption{Maximum pushover capacity for the seven building configurations.}
\label{tab:ch4_peak_capacity}
\centering
\footnotesize
\setlength{\tabcolsep}{3pt}
\begin{tabular}{@{}l c c c r r@{}}
\toprule
\textbf{Model} & \textbf{Soil} & \textbf{Z} & \textbf{Dir.} & \shortstack{\textbf{Peak base shear}\\\textbf{(kN)}} & \shortstack{\textbf{Displacement at peak}\\\textbf{(mm)}}\\
\midrule
Regular & C & 0.40 & X & 1836.338 & 198.040\\
L-shaped & C & 0.40 & Y & 1800.120 & 187.539\\
L1-shaped & C & 0.40 & X & 1606.250 & 212.287\\
L2-shaped & C & 0.35 & Y & 1804.152 & 172.789\\
T-shaped & C & 0.40 & Y & 1640.715 & 223.937\\
U-shaped & C & 0.40 & Y & 1882.070 & 174.697\\
Non-Parallel & C & 0.40 & Y & 1827.694 & 154.530\\
\bottomrule
\end{tabular}
\end{table}

Figure~\ref{fig:nbc105_2025_capacity_soil_C_Z_0p4_x} and Figure~\ref{fig:nbc105_2025_capacity_soil_C_Z_0p4_y} show the Soil Type C, $Z=0.40$ capacity curves in the X and Y directions through the first peak. Performance points are discussed separately.

\begin{figure}[!htbp]
\centering
\includegraphics[width=0.90\textwidth]{figures/ch4/authority_exact/fig_4_1_CP97_CP66_exact_clip.pdf}
\caption{Capacity curves in the X direction.}
\label{fig:nbc105_2025_capacity_soil_C_Z_0p4_x}
\end{figure}

\begin{figure}[!htbp]
\centering
\includegraphics[width=0.90\textwidth]{figures/ch4/authority_exact/fig_4_2_CP97_CP66_exact_clip.pdf}
\caption{Capacity curves in the Y direction.}
\label{fig:nbc105_2025_capacity_soil_C_Z_0p4_y}
\end{figure}

\begin{table}[!htbp]
\caption[Maximum performance-point shear and displacement (seven models)]{Maximum performance-\allowbreak{}point base shear and roof displacement for the seven-model comparison.}
\label{tab:ch4_pp_summary_revised}
\centering
\footnotesize
\setlength{\tabcolsep}{3pt}
\begin{tabular}{@{}l r c c c r c c c@{}}
\toprule
\textbf{Model} & \shortstack{\textbf{Max. PP base}\\\textbf{shear (kN)}} & \textbf{Soil} & \textbf{Z} & \textbf{Dir.} & \shortstack{\textbf{Max. PP roof}\\\textbf{disp. (mm)}} & \textbf{Soil} & \textbf{Z} & \textbf{Dir.}\\
\midrule
Regular & 456.253 & C & 0.40 & X & 26.053 & C & 0.40 & X\\
L-shaped & 407.295 & C & 0.40 & X & 25.900 & C & 0.40 & X\\
L1-shaped & 340.753 & C & 0.40 & X & 24.592 & C & 0.40 & X\\
L2-shaped & 370.935 & C & 0.40 & X & 21.171 & C & 0.40 & X\\
T-shaped & 343.488 & C & 0.40 & X & 24.802 & C & 0.40 & X\\
U-shaped & 435.084 & C & 0.40 & X & 27.072 & C & 0.40 & X\\
Non-Parallel & 399.636 & C & 0.40 & X & 21.328 & C & 0.40 & X\\
\bottomrule
\end{tabular}
\end{table}

Figure~\ref{fig:nbc105_2025_storey_shear_soil_C_Z_0p4}, Figure~\ref{fig:nbc105_2025_drift_soil_C_Z_0p4}, and Figure~\ref{fig:nbc105_2025_cmcr_main_story} summarize the storey-shear, ULS drift, and main-storey CM--CR separation results. The maximum ULS drift ratio remains below the NBC 105:2025 ULS drift limit. CM--CR separation is interpreted descriptively and is not used as a formal torsional-irregularity classification.

\begin{figure}[!htbp]
\centering
\includegraphics[width=0.96\textwidth]{figures/ch4/comment31/story_shear_levels_source_backed.pdf}
\caption{Storey shear profiles for Soil Type C, $Z=0.40$; recorded EQx values by elevation.}
\label{fig:nbc105_2025_storey_shear_soil_C_Z_0p4}
\end{figure}

\begin{figure}[!htbp]
\centering
\includegraphics[width=0.96\textwidth]{figures/ch4/comment31/uls_drift_profile_X_source_backed.pdf}
\caption{ULS drift profiles in the X direction for Soil Type C, $Z=0.40$.}
\label{fig:nbc105_2025_drift_soil_C_Z_0p4}
\end{figure}

\begin{figure}[!htbp]
\centering
\includegraphics[width=0.90\textwidth]{figures/ch4/authority_exact/fig_4_5_CP110_exact.pdf}
\caption{Maximum CM--CR separation.}
\label{fig:nbc105_2025_cmcr_main_story}
\end{figure}

\begin{table}[!htbp]
\caption{Maximum governing ULS drift ratios for Soil Type C, $Z=0.40$.}
\label{tab:ch4_drift_summary}
\centering
\footnotesize
\setlength{\tabcolsep}{3pt}
\begin{tabular}{@{}l c c c c c c@{}}
\toprule
\textbf{Model} & \textbf{Soil} & \textbf{Z} & \textbf{Direction} & \shortstack{\textbf{Governing}\\\textbf{Story}} & \shortstack{\textbf{Max ULS}\\\textbf{Drift Ratio}} & \textbf{Utilization}\\
\midrule
Regular & C & 0.400 & X & Story2 & 0.004784 & 0.191\\
L & C & 0.400 & X & Story2 & 0.004758 & 0.190\\
L1 & C & 0.400 & X & Story2 & 0.004570 & 0.183\\
L2 & C & 0.400 & X & Story2 & 0.004025 & 0.161\\
T & C & 0.400 & Y & Story2 & 0.004618 & 0.185\\
U & C & 0.400 & X & Story2 & 0.004954 & 0.198\\
Non-Parallel & C & 0.400 & X & Story2 & 0.004041 & 0.162\\
\bottomrule
\end{tabular}
\end{table}

\begin{table}[!htbp]
\caption{Maximum resultant CM--CR separation across main Stories 1--3.}
\label{tab:ch4_cmcr_summary}
\centering
\footnotesize
\setlength{\tabcolsep}{3pt}
\begin{tabular}{@{}l c c c r@{}}
\toprule
\textbf{Model} & \textbf{Soil} & \textbf{Z} & \shortstack{\textbf{Storey at}\\\textbf{maximum}} & \shortstack{\textbf{Maximum resultant}\\\textbf{separation (m)}}\\
\midrule
Regular & D & 0.350 & Story2 & 0.398\\
L & C & 0.250 & Story1 & 0.581\\
L1 & C & 0.250 & Story1 & 0.640\\
L2 & C & 0.250 & Story1 & 1.186\\
T & C & 0.300 & Story2 & 0.475\\
U & C & 0.250 & Story2 & 0.424\\
NP & C & 0.250 & Story2 & 1.237\\
\bottomrule
\end{tabular}
\end{table}


\subsection{Complete Governing Drift Profiles}
Figures~\ref{fig:drift_profiles_cp112a} and~\ref{fig:drift_profiles_cp112b} expand the governing-drift comparison from a single maximum to the complete main-storey ULS profiles for Soil Type C, $Z=0.40$. The profiles confirm that Story~2 governs the reported maximum for every model; the T configuration is the only model in which the Y-direction value slightly exceeds the X-direction value at the governing storey (0.004618 versus 0.004597).

\begin{figure}[!htbp]
\centering
\begin{subfigure}[t]{0.48\textwidth}\centering
\includegraphics[width=\linewidth]{figures/ch4/cp112_generated/drift_regular.png}\subcaption{Regular}
\end{subfigure}\hfill
\begin{subfigure}[t]{0.48\textwidth}\centering
\includegraphics[width=\linewidth]{figures/ch4/cp112_generated/drift_l.png}\subcaption{L-shaped}
\end{subfigure}

\begin{subfigure}[t]{0.48\textwidth}\centering
\includegraphics[width=\linewidth]{figures/ch4/cp112_generated/drift_l1.png}\subcaption{L1-shaped}
\end{subfigure}\hfill
\begin{subfigure}[t]{0.48\textwidth}\centering
\includegraphics[width=\linewidth]{figures/ch4/cp112_generated/drift_l2.png}\subcaption{L2-shaped}
\end{subfigure}
\caption{ULS drift profiles: Regular, L, L1 and L2.}
\label{fig:drift_profiles_cp112a}
\end{figure}

Figure~\ref{fig:drift_profiles_cp112a} shows that the governing drift occurs at Story~2 for the Regular, L, L1 and L2 configurations under the matched Soil Type C, $Z=0.40$ case. The profile shapes are similar in general form, but the magnitudes differ, with L2 showing the lowest governing drift among these four models.


\begin{figure}[!htbp]
\centering
\begin{subfigure}[t]{0.48\textwidth}\centering
\includegraphics[width=\linewidth]{figures/ch4/cp112_generated/drift_t.png}\subcaption{T-shaped}
\end{subfigure}\hfill
\begin{subfigure}[t]{0.48\textwidth}\centering
\includegraphics[width=\linewidth]{figures/ch4/cp112_generated/drift_u.png}\subcaption{U-shaped}
\end{subfigure}

\begin{subfigure}[t]{0.48\textwidth}\centering
\includegraphics[width=\linewidth]{figures/ch4/cp112_generated/drift_non_parallel.png}\subcaption{Non-Parallel}
\end{subfigure}
\caption{ULS drift profiles: T, U and Non-Parallel.}
\label{fig:drift_profiles_cp112b}
\end{figure}

Figure~\ref{fig:drift_profiles_cp112b} confirms Story~2 as the governing level for the T, U and Non-Parallel configurations. The U-shaped model has the largest governing drift in the seven-model comparison, while the Non-Parallel model remains lower under the same matched demand case.

\FloatBarrier

\subsection{Percentage Change Relative to the Regular Reference}
Table~\ref{tab:percent_change_regular_cp112} expresses selected response quantities relative to the Regular reference using $(R_i-R_{reg})/R_{reg}\times100$. The storey-shear percentage is based on the available EQx Soil C, $Z=0.40$ Story-1 results because a corresponding EQy shear series is unavailable.

\begin{table}[!htbp]
\centering
\caption{Percentage change from the Regular reference for selected response quantities.}
\label{tab:percent_change_regular_cp112}
\scriptsize
\resizebox{\textwidth}{!}{%
\begin{tabular}{@{}lrrrrr@{}}
\toprule
Model & EQx shear & Max ULS drift & Max CM--CR & Peak pushover shear & Max PP disp. \\
 & (\%) & (\%) & (\%) & (\%) & (\%) \\
\midrule
Regular & 0.0 & 0.0 & 0.0 & 0.0 & 0.0 \\
L & -10.6 & -0.5 & +46.0 & -2.0 & -0.6 \\
L1 & -24.2 & -4.5 & +60.8 & -12.5 & -5.6 \\
L2 & -14.2 & -15.9 & +198.0 & -1.8 & -18.7 \\
T & -23.5 & -3.5 & +19.3 & -10.7 & -4.8 \\
U & -5.3 & +3.6 & +6.5 & +2.5 & +3.9 \\
Non-Parallel & -8.6 & -15.5 & +210.8 & -0.5 & -18.1 \\
\bottomrule
\end{tabular}
}
\end{table}

The comparison demonstrates why a single irregularity ranking is inappropriate. L2 and Non-Parallel show much larger CM--CR separation than the Regular reference while simultaneously exhibiting lower governing drift and lower maximum performance-point displacement. Conversely, U is close to the Regular model in EQx storey shear but exceeds it slightly in governing drift, first-peak pushover capacity, and maximum performance-point displacement. Geometry therefore changes different response measures in different directions.

\begin{table}[!htbp]
\centering
\caption{Detailed CM--CR coordinates at the governing main-storey separation for each model.}
\label{tab:cmcr_detailed_cp112}
\scriptsize
\resizebox{\textwidth}{!}{%
\begin{tabular}{@{}llrrrrrrr@{}}
\toprule
Model & Case / story & CM-X & CM-Y & CR-X & CR-Y & $\Delta X$ & $\Delta Y$ & Resultant \\
 & & \multicolumn{7}{c}{(m)} \\
\midrule
Regular & D,0.35 / Story2 & 5.6885 & 3.1168 & 6.0619 & 2.9802 & -0.3734 & 0.1366 & 0.3976 \\
L & C,0.25 / Story1 & 4.8491 & 3.3662 & 5.4237 & 3.2835 & -0.5746 & 0.0827 & 0.5805 \\
L1 & C,0.25 / Story1 & 4.5223 & 3.8090 & 5.1408 & 3.6464 & -0.6185 & 0.1626 & 0.6395 \\
L2 & C,0.25 / Story1 & 3.7215 & 3.1287 & 4.8085 & 3.6021 & -1.0870 & -0.4734 & 1.1856 \\
T & C,0.30 / Story2 & 5.6314 & 3.7622 & 6.1018 & 3.6945 & -0.4704 & 0.0677 & 0.4752 \\
U & C,0.25 / Story2 & 5.6880 & 2.9827 & 6.0731 & 2.8046 & -0.3851 & 0.1781 & 0.4243 \\
Non-Parallel & C,0.25 / Story2 & 3.7987 & 3.1802 & 4.9046 & 3.7352 & -1.1059 & -0.5550 & 1.2374 \\
\bottomrule
\end{tabular}
}
\end{table}

The large separations in L2 and Non-Parallel arise from different geometric mechanisms. L2 combines a stepped re-entrant topology with an asymmetric distribution of plan area, while the Non-Parallel model combines area deficit with skewed framing orientation. The Regular, L, L1, T and U cases have smaller governing separations, but their maxima occur under different demand cases and storeys. CM--CR separation is therefore retained as a descriptive mass--stiffness eccentricity measure, not as a substitute for the NBC diaphragm-edge torsional-irregularity ratio.
\FloatBarrier



\section{Phase IV: Determination of NBC 205 RUD Applicability Limits}
\subsection{Overall Comparative Seismic Performance}
The results show that different building configurations become critical for different response parameters. In the storey-shear results, the Regular building governs the maximum lower-storey shear. The U-shaped building governs the first-peak pushover base shear and the largest performance-\allowbreak{}point roof displacement, while the Regular building governs the highest performance-\allowbreak{}point base shear. The Non-Parallel building governs the main-storey CM--CR separation comparison.

This distribution of critical responses demonstrates that the controlling model depends on the response quantity being examined. The Regular model governs the storey-shear demand and performance-\allowbreak{}point base shear, whereas the U-shaped model governs the first-peak pushover capacity and performance-\allowbreak{}point displacement. The irregular configurations must therefore be interpreted from their full response patterns rather than from a single ranking.

The effect of seismic zone factor was clear in the storey-shear, capacity, and performance-point results: their governing values occurred under Soil Type C with $Z=0.40$. Soil Type D at $Z=0.35$ was included in the comparison but did not govern these quantities. This does not imply that soil effects are unimportant; rather, it reflects the controlling combinations present in the analysed cases.

Overall seismic performance cannot be judged from a single maximum value or from a simple ranking without a formal weighting method. High ultimate pushover strength must be considered together with drift, CM--CR separation, performance-point response, final-step hinge evidence, and the available design-check observations. The seven-model comparison highlights how different plan configurations respond to seismic demand, but does not establish general code acceptance for RUD detailing in irregular buildings.

The Regular NBC 205:2024 model is retained as the study reference benchmark. The results do not justify automatic use of RUD detailing for irregular buildings. The applicability of RUD detailing to any of the six plan-irregular configurations must be evaluated for the particular building rather than inferred from a single response ranking. The separate four-storey extension also requires building-specific assessment; its design-check findings should not be interpreted as a universal storey limit.

\subsection{Geometric Irregularity Index (GII)}
The Geometric Irregularity Index (GII) was examined to assess whether plan-area irregularity is consistently related to the observed performance-point response. In this study, GII represents the percentage reduction in plan area relative to the corresponding bounding rectangle. The GII values range from 0\% for the Regular model to 33.33\%, while the Non-Parallel model also has a non-parallelism angle of $\theta = 36.87^{\circ}$.

\begin{table}[!htbp]
\centering
\caption[GII and seismic-response summary for seven configurations]{Geometric Irregularity Index (GII) and response summary for all seven configurations.}
\label{tab:gii_response_summary}
\small
\resizebox{\textwidth}{!}{%
\begin{tabular}{@{}lccccc@{}}
\toprule
Model & GII (\%) & $T_1$ (s) & Peak base shear (kN) & Max ULS drift ratio & Max PP shear / disp. \\
\midrule
Regular & 0.00 & 0.699 & 1836.338 & 0.004784 & 456.253 / 26.053 \\
L & 16.67 & 0.698 & 1800.120 & 0.004758 & 407.295 / 25.900 \\
L1 & 33.33 & 0.685 & 1606.250 & 0.004570 & 340.753 / 24.592 \\
L2 & 33.33 & 0.649 & 1804.152 & 0.004025 & 370.935 / 21.171 \\
T & 33.33 & 0.698 & 1640.715 & 0.004618 & 343.488 / 24.802 \\
U & 16.67 & 0.707 & 1882.070 & 0.004954 & 435.084 / 27.072 \\
Non-Parallel & 22.22 & 0.650 & 1827.694 & 0.004041 & 399.636 / 21.328 \\
\bottomrule
\end{tabular}
}
\par\smallskip
\begin{minipage}{0.96\textwidth}\footnotesize
$T_1$ is the Mode-1 period for Soil Type C, $Z=0.40$. Peak base shear is the maximum first-global-peak pushover value. Performance point is maximum base shear (kN) / maximum roof displacement (mm) from the available nonlinear-analysis results.
\end{minipage}
\end{table}

\begin{figure}[!htbp]
\centering
\includegraphics[width=0.94\textwidth]{figures/ch4/comment31/gii_pp_base_shear_distinct_markers.pdf}
\caption{Geometric Irregularity Index (GII) and performance-point base shear.}
\label{fig:gii_pp_base_shear}
\end{figure}

\begin{figure}[!htbp]
\centering
\includegraphics[width=0.94\textwidth]{figures/ch4/comment31/gii_pp_roof_displacement_distinct_markers.pdf}
\caption{Geometric Irregularity Index (GII) and performance-point roof displacement.}
\label{fig:gii_pp_roof_disp}
\end{figure}
\FloatBarrier

The results do not show a clear monotonic relationship between GII and performance-point response. For example, the L- and U-shaped models have the same GII of 16.67\%, but their performance-point base shear and roof displacement are 407.295 kN and 25.900 mm, and 435.084 kN and 27.072 mm, respectively. Similarly, the L1-, L2-, and T-shaped models all have GII = 33.33\%, but their performance-point responses differ considerably.

The U-shaped model further illustrates that a relatively low GII does not necessarily correspond to lower response, as it exhibited the largest performance-point roof displacement (27.072 mm) and the highest first-peak pushover base shear (1882.070 kN). The Non-Parallel model also demonstrates that area-based GII does not capture other important geometric characteristics such as framing orientation and torsional response.

Therefore, for the seven configurations investigated, GII alone was insufficient to explain the variation in seismic response. It remains useful as a geometric descriptor of plan-area irregularity but should not be used independently as a seismic performance criterion or as the sole basis for assessing the applicability of RUD provisions. The effect of irregularity should be evaluated together with drift, CM--CR eccentricity, torsional rotation, pushover capacity, performance-point response, and hinge behaviour. Further investigation using a larger range of plan configurations would be required to establish a statistically supported relationship between GII and seismic response.



\begin{figure}[!htbp]
\centering
\begin{subfigure}[t]{0.48\textwidth}\centering
\includegraphics[width=\linewidth]{figures/ch4/cp112_generated/gii_max_drift.png}\subcaption{ULS drift}
\end{subfigure}\hfill
\begin{subfigure}[t]{0.48\textwidth}\centering
\includegraphics[width=\linewidth]{figures/ch4/cp112_generated/gii_cmcr.png}\subcaption{CM--CR separation}
\end{subfigure}

\begin{subfigure}[t]{0.48\textwidth}\centering
\includegraphics[width=\linewidth]{figures/ch4/cp112_generated/gii_peak_pushover.png}\subcaption{Peak pushover shear}
\end{subfigure}\hfill
\begin{subfigure}[t]{0.48\textwidth}\centering
\includegraphics[width=\linewidth]{figures/ch4/cp112_generated/gii_pp_displacement.png}\subcaption{Performance-point roof displacement}
\end{subfigure}
\caption{Geometric Irregularity Index (GII) and selected seismic responses.}
\label{fig:gii_multiresponse_cp112a}
\end{figure}

Figure~\ref{fig:gii_multiresponse_cp112a} demonstrates that the same or similar GII can correspond to different drift, CM--CR separation, peak pushover strength and performance-point displacement. The plotted responses therefore do not support treating GII as a unique predictor of seismic behaviour.


The performance-point-shear comparison is already presented in Figure~\ref{fig:gii_pp_base_shear}. Together with the four-response comparison in Figure~\ref{fig:gii_multiresponse_cp112a}, these results indicate that plan topology and stiffness distribution cannot be described by area deficit alone.

\FloatBarrier

\begin{table}[!htbp]
\centering
\caption{Final seven-model comparative summary.}
\label{tab:final_summary_cp112}
\scriptsize
\resizebox{\textwidth}{!}{%
\begin{tabular}{@{}lrrrrrrp{2.7cm}@{}}
\toprule
Model & GII & Max drift & Max CM--CR & EQx shear & Peak push. & Max PP disp. & Hinge evidence \\
 & (\%) & ratio & (m) & (kN) & (kN) & (mm) & \\
\midrule
Regular & 0.00 & 0.004784 & 0.398 & 532.489 & 1836.338 & 26.053 & Final-step X/Y images; member-level table unavailable \\
L & 16.67 & 0.004758 & 0.581 & 476.287 & 1800.120 & 25.900 & Same evidence boundary \\
L1 & 33.33 & 0.004570 & 0.640 & 403.859 & 1606.250 & 24.592 & Same evidence boundary \\
L2 & 33.33 & 0.004025 & 1.186 & 456.660 & 1804.152 & 21.171 & Same evidence boundary \\
T & 33.33 & 0.004618 & 0.475 & 407.396 & 1640.715 & 24.802 & Same evidence boundary \\
U & 16.67 & 0.004954 & 0.424 & 504.048 & 1882.070 & 27.072 & Same evidence boundary \\
Non-Parallel & 22.22 & 0.004041 & 1.237 & 486.757 & 1827.694 & 21.328 & Same evidence boundary \\
\bottomrule
\end{tabular}
}
\par\smallskip
\begin{minipage}{0.97\textwidth}\footnotesize
EQx shear is the validated Soil C, $Z=0.40$ Story-1 value. Formal NBC torsional-irregularity classification is not assigned from CM--CR separation. The table summarizes the evidence used in this thesis; it is not a code-acceptance table.
\end{minipage}
\end{table}

\clearpage
\subsection{Alignment of Research Questions, Results and Conclusions}
The six research objectives are linked to their corresponding Chapter~4 evidence and Chapter~5 findings in Table~\ref{tab:rq_objective_traceability_cp112}. This mapping provides a guide to the manuscript argument; it does not constitute an additional dataset.
\begin{table}[!ht]
\centering
\caption{Alignment of the six research questions/objectives with results and conclusions.}
\label{tab:rq_objective_traceability_cp112}
\scriptsize
\begin{tabular}{@{}c p{8.0cm} p{3.4cm}@{}}
\toprule
No. & Result evidence in Chapter 4 & Conclusion linkage in Chapter 5 \\
\midrule
1 & Phase I concrete design-check evidence for the 2012 reference & 2012 design-check finding \\
2 & Regular benchmark modal/storey/drift/pushover/performance response & Regular benchmark finding \\
3 & Seven-model plan-irregular comparison and complete drift/force/capacity tables & Irregular-model finding \\
4 & CM--CR coordinates, drift profiles, pushover and performance-point comparisons & CM--CR/response finding \\
5 & Phase IV applicability-boundary synthesis and final comparative table & RUD scope/height finding \\
6 & Seven-model GII tables and selected GII--response comparisons & GII response finding \\
\bottomrule
\end{tabular}
\end{table}
\FloatBarrier

